\section{Greedy deletion in a fitted boosted forest}
\label{sec:tree-minpi}
Let $N$ be the $p$ input columns of a fitted forest. For a retained set
$S\subseteq N$, traverse each tree until a leaf or the first split on a column
outside $S$. At the latter node return the mean of the two child margins
weighted by their training Hessian cover, applying the same definition
recursively down to the leaves. Native missing-value routing still applies to
retained columns. Let $f_S$ collect the resulting raw ensemble margins on the
calibration rows, with one coordinate per row and output. The fidelity budget is
$B=\rho^2\|f_N-f_\varnothing\|^2$, with $\rho=.20$ in the primary experiment.
A deletion is feasible when its squared discrepancy from $f_N$ is at most $B$.
This criterion uses the full predictor as the reference and uses no calibration
outcomes.

We precompute the child-minus-parent margin change for each edge traversed by
each calibration row. Such changes telescope along tree paths, as in the
feature contribution method of \citet{palczewska2013forest}. An edge in our
truncated predictor survives only if the retained set contains \emph{every}
split feature leading to it, including its parent's split. The joint
contribution routine in \texttt{treeinterpreter}~\cite{saabasTreeInterpreter}
already groups node increments by this set of prefix features. Use $a_H$
for the resulting vector of increments with required set $H$.
Then $f_S=f_\varnothing+\sum_{H\subseteq S}a_H$. At deployment the trees can
be traversed with the retained raw columns alone.

Write $r=f_N-f_S$ and let $\delta_j=\sum_{j\in H\subseteq S}a_H$ be the change
from deleting $j\in S$. Greedy deletion removes the feasible column with the
smallest resulting squared error,
\begin{equation}
 L_j=\|r+\delta_j\|^2=c_j+\|\delta_j\|^2,\qquad
 c_j=\|r\|^2+2\langle r,\delta_j\rangle .
 \label{eq:tree-error}
\end{equation}
It stops when every remaining deletion exceeds $B$. We retain the exact cross
term by summing $\langle r,a_H\rangle$ over active groups requiring $j$. The
expensive operation is assembling $\delta_j$ across calibration rows to
recalculate its squared norm after another column has been removed.

\paragraph{Votes by coalition size.}
For each feature $j$ and size $s$, take a uniformly chosen $s$-subset $T$
of $N\setminus\{j\}$. The direction $h_{js}$ is the expected change in raw
predictions when $j$ joins $T$. Each path group is a unanimity game on its
required features~\citep{shapley1953value}. If $j\in H$, adding $j$ activates
that group precisely when $T$ contains the other $q_H=|H|-1$ features of $H$.
Counting subsets gives the exact direction
\begin{equation}
 h_{js}=\mathbb E[f_{T\cup\{j\}}-f_T\mid |T|=s]
       =\sum_{H\ni j}\frac{\binom{s}{q_H}}{\binom{p-1}{q_H}}\,a_H .
 \label{eq:tree-size}
\end{equation}
A numerator is zero for $s<q_H$. Small coalitions give greater relative weight
to increments that require fewer distinct columns. At $s=p-1$ every coefficient is one and
$h_{j,p-1}=f_N-f_{N\setminus\{j\}}$. Size-conditioned marginals are the
ingredients of symmetric semivalues~\citep{dubey1981semivalues}.

Use the distinct sizes in $\mathcal S=\{0,\lfloor(p-1)/2\rfloor,p-1\}$.
At a cutoff $b$, the size-$s$ vote passes if $c_j+Q_{js}\le b$, where
\begin{equation}
 Q_{js}=\frac{\langle h_{js},\delta_j\rangle^2}{\|h_{js}\|^2},\qquad
 \min_{s\in\mathcal S}(b-c_j-Q_{js})=b-c_j-\max_{s\in\mathcal S}Q_{js}.
 \label{eq:tree-vote}
\end{equation}
Set $Q_{js}=0$ when $h_{js}=0$. Each vote measures the squared projection onto
one line, admitting both orientations of that direction. A failed vote rules
out the current deletion. The original MinPI takes a minimum of changes in a
utility-threshold vote across individual coalitions~\citep{shi2026minpi}. Here
each size supplies a marginal \emph{prediction vector}; the minimum acts on
its deletion-feasibility vote.

For example, take a scalar game with $a_{\{1\}}=a_{\{2\}}=1$ and
$a_{\{1,2\}}=-1$. Removing feature 2 from the full model leaves its prediction
at one. Removing feature 1 next would reduce it to zero. Its full-size direction
$h_{1,1}$ vanishes, while $h_{1,0}=1$ gives the exact squared loss of one.
Only the small-coalition vote vetoes the second deletion when $0\le b<1$.

\paragraph{Full-coalition controls.}
A one-direction control uses $s=p-1$ alone. A second construction splits the
coordinates of $h_{j,p-1}$ by sign ($>0$) and by whether absolute magnitude
exceeds its root mean square. Let $v_b$ be this full-coalition vector masked
to cell $b$, and $V_j$ the span of the nonzero cell vectors. Their disjoint
supports give
\[
 Q_j=\sum_{v_b\ne0}\frac{\langle v_b,\delta_j\rangle^2}{\|v_b\|^2}
     = \max_{u\in V_j,\,\|u\|\le1}\langle u,\delta_j\rangle^2 .
\]
This sum evaluates the minimum slack over all linear combinations of the four
directions. Different coalition sizes in~\eqref{eq:tree-vote} cast separate
votes.

\begin{proposition}
At each deletion step, testing the projection votes and evaluating the surviving
candidates exactly gives the same deletion, including the chosen tie, as
evaluating~\eqref{eq:tree-error} for every candidate in exact arithmetic.
\end{proposition}
\begin{proof}
For any one of the tested spaces $V$ with projection $P_V$, the gap between
$L_j$ and $c_j+\|P_V\delta_j\|^2$ is $\|\delta_j-P_V\delta_j\|^2\ge0$
by the orthogonal projection identity~\citep{melrose2013bessel}. Thus the
largest of the tested lower bounds is still valid. Begin with cutoff $B$ and lower it
whenever a better feasible candidate is evaluated exactly. A candidate whose
lower bound strictly exceeds the cutoff cannot win. Candidates at equality remain
eligible. Equal errors are resolved by removing the later member of the
descending Gain order, which ranks larger column IDs first at tied Gain.
Applying the argument after each deletion gives identical paths.
\end{proof}

Initially $\delta_j=h_{j,p-1}$, so the full-size bound is exact. After each deletion only groups
that lose a required column are removed, and their contributions to
the tested projections are subtracted from cached inner products.
Previously computed exact norms remain valid for unaffected features. An
unresolved feature receives a fresh exact norm only if its lower bound can
still win. Maintaining these bounds incurs an initialization and update cost.
Greedy termination certifies that no further
single deletion meets the calibration budget at the returned support.
